% !TeX root = ../imdt-proposal-report.tex
% 这是唯一的日常配置入口。把值填在花括号内，两份材料自动共用。
\imdtsetup{
  mode        = writing,      % writing 正式写作；reference 查看原 Word 的完整 12 页
  hints       = show,         % show 显示填写提示；hide 隐藏红色/灰色提示
  examples    = show,         % show 显示示例图表公式和演示引文；hide 全部隐藏
  fonts       = auto,         % 私人 fonts/ 有原字库时自动使用；portable 仅供预览
  output      = electronic,  % electronic 电子稿；print 双面打印
  title       = {},          % 论文题目，不超过 25 个汉字（符）
  author      = {},          % 你的姓名
  student-id  = {},          % 你的学号，例如：2026123456
  school      = {深圳国际研究生院},
  major       = {电子信息（互动媒体设计与技术）},
  supervisor  = {},          % 学籍系统登记的导师
  supervisor-rank = {},      % 计划表导师职称，例如：教授
  enrollment-date = {},      % 计划表入学年月，例如：2026年9月
  report-date = {\hspace{1em}年\hspace{1em}月\hspace{1em}日},
  form-date   = {年\hspace{4.5em}月\hspace{5em}日},
  defense-date = {},         % 例如：2028年6月
  bibliography = {bib/references},
  chapters = {
    chapters/01-background,
    chapters/02-related-work,
    chapters/03-research-plan,
    chapters/04-outcomes-schedule
  },
}

% 计划表专属配置：用 yes/no 或指定名称选择，不必自己定位或画勾。
% 没有填写的选择项默认留空；不需要的条目保留注释。
% \imdtsetup{
%   topic-scope = {选题范围}, meeting-date = {2026年10月21日},
%   secrecy = public,                 % public / secret / confidential / blank
%   cross-discipline-check = no,       % yes / no / blank
%   expert-review = no,                % yes / no / blank
%   experiment = yes, safety-training = yes, safety-report = yes,
%   thesis-type = product,            % research / survey / case / product / scheme / other
%   domestic-references = {12}, foreign-references = {10},
%   research-conditions = {在此写已有设备与研究条件。},
%   project-category = {在此写科研项目情况。},
%   work-plan = {第一阶段完成文献调研。\par 第二阶段实现并评价原型。},
%   literature-due = {2026年12月}, research-due = {2027年9月},
%   writing-due = {2028年3月},
% }

% 提交时取消下一行注释：一次关闭全部提示、演示内容与示例附录。
% \imdtsubmission

% 大稿件：先完整编译，再按需取消注释以只编译某章；提交前重新注释。
% \includeonly{chapters/03-research-plan}
